\section{The SMC Algorithm}
\label{sec:algorithm}

\subsection{Setup and Notation}

Let $G = (V, E)$ be the census-tract adjacency graph for a state, with
$|V| = n$ tracts and $k$ target districts.
A redistricting plan $\pi : V \to [k]$ is \emph{valid} if every district is
connected and population-balanced within tolerance $\epsilon$:
\[
  \bigl|\mathrm{pop}(\pi^{-1}(d)) - P_{\mathrm{ideal}}\bigr|
  \;\leq\; \epsilon \cdot P_{\mathrm{ideal}}
  \quad\text{for all } d \in [k],
\]
where $P_{\mathrm{ideal}} = \sum_v \mathrm{pop}(v) / k$ and $\epsilon = 0.005$
(half a percent).
Let $\Pi$ denote the set of all valid plans.
The target distribution is the uniform distribution $\pi^{\mathrm{unif}}$ over
$\Pi$.

An \smc{} particle at stage $t$ is a \emph{partial plan}: an assignment
$\pi^{(t)} : V \to [t] \cup \{\mathrm{unassigned}\}$, where districts
$1, \ldots, t$ are finalised and the remaining $n - |\pi^{-1}([t])|$ tracts
are unassigned.
The \emph{unassigned subgraph} $G^{(t)} = G[V \setminus \pi^{-1}([t])]$ is
the induced subgraph on unassigned tracts.

\subsection{Deterministic Seed Schedule}

The \smc{} implementation uses two deterministic seed derivations to ensure
full reproducibility from a single \texttt{base\_seed}.

\paragraph{Per-particle seed.}
Particle $i$ at stage $t$ uses the seed:
\[
  \mathrm{particle\_seed}(s, t, i)
  = \mathrm{SHA256}\!\left(
    \texttt{"SMC\_PARTICLE\_"} \,\|\, t^{\mathrm{le32}}
    \,\|\, \texttt{"\_"} \,\|\, i^{\mathrm{le32}}
    \,\|\, \texttt{"\_"} \,\|\, s^{\mathrm{le64}}
  \right)[0{:}8] \;\text{as } u64,
\]
where $s$ is the \texttt{base\_seed}, $\cdot^{\mathrm{le32}}$ denotes
little-endian 4-byte encoding, $\cdot^{\mathrm{le64}}$ denotes little-endian
8-byte encoding, and $[0{:}8]$ takes the first 8 bytes of the SHA-256 digest.

\paragraph{Resampling seed.}
The systematic resampling at round $r$ uses:
\[
  \mathrm{resample\_seed}(s, r)
  = \mathrm{SHA256}\!\left(
    \texttt{"SMC\_RESAMPLE\_"} \,\|\, r^{\mathrm{le32}}
    \,\|\, \texttt{"\_"} \,\|\, s^{\mathrm{le64}}
  \right)[0{:}8] \;\text{as } u64.
\]

These derivations ensure that all randomness is determined by \texttt{base\_seed}
alone: any verifier who knows \texttt{base\_seed} can reconstruct every
per-particle and per-resampling RNG state without additional secrets.

\subsection{The Full SMC Algorithm}

\begin{algorithm}
\caption{\smc{} for Redistricting}
\label{alg:smc}
\begin{algorithmic}[1]
\Require Adjacency graph $G$, population vector $\mathrm{pop}$, district
         count $k$, particle count $N$, tolerance $\epsilon$,
         resample threshold $\tau$ (default $0.5$), base seed $s$
\Ensure  $(N$ weighted plans$,\, \mathbf{w})$ with $\sum_i w_i = 1$
\State $\mathbf{p}[i] \gets \mathrm{empty\_partial\_plan}$ for $i = 1, \ldots, N$
\State $\log \mathbf{w}[i] \gets 0$ for $i = 1, \ldots, N$
\State $\mathrm{round} \gets 0$
\For{$\mathrm{stage} \gets 1$ \textbf{to} $k - 1$}
  \ParFor{$i \gets 1$ \textbf{to} $N$} \Comment{Rayon \texttt{par\_iter()}}
    \State $\mathrm{rng} \gets \mathrm{SmallRng::seed\_from\_u64}(\mathrm{particle\_seed}(s, \mathrm{stage}, i))$
    \State $C \gets \mathrm{largest\_connected\_component}(G^{(\mathrm{stage}-1)}[\mathbf{p}[i]])$
    \State $v^* \gets \mathrm{rng.choose}(C)$ \Comment{uniform seed in $C$}
    \State $T \gets \mathrm{Wilson\_UST}(G[C], \mathrm{rng})$ \Comment{uniform spanning tree}
    \State $\mathcal{V} \gets \{e \in T : \text{removing } e \text{ gives a balanced bipartition of } C\}$
    \If{$\mathcal{V} = \emptyset$}
      \State $\log \mathbf{w}[i] \gets -\infty$ \Comment{kill particle}
    \Else
      \State $e^* \gets \mathrm{rng.choose}(\mathcal{V})$ \Comment{uniform cut edge}
      \State assign the component of $C$ containing $v^*$ after removing $e^*$ as district $\mathrm{stage}$
      \State $\log \mathbf{w}[i] \mathrel{+}= \log|\mathcal{V}|$ \Comment{importance weight increment}
    \EndIf
  \EndParFor
  \State $\ess \gets \bigl(\sum_i \exp(\log \mathbf{w}[i])\bigr)^2 \;\big/\; \sum_i \exp(2 \log \mathbf{w}[i])$
  \If{$\ess < \tau \cdot N$}
    \State $(\mathbf{p}, \mathrm{index\_map}) \gets \mathrm{systematic\_resample}(\mathbf{p}, \log \mathbf{w},\, \mathrm{resample\_seed}(s, \mathrm{round}))$
    \State $\log \mathbf{w}[i] \gets 0$ for all $i$
    \State $\mathrm{round} \mathrel{+}= 1$
  \EndIf
\EndFor
\State assign all remaining unassigned tracts as district $k$ \Comment{for each particle}
\State $\mathbf{w} \gets \mathrm{kahan\_softmax}(\log \mathbf{w})$ \Comment{normalised weights}
\State \Return $(\mathbf{p}, \mathbf{w})$
\end{algorithmic}
\end{algorithm}

The \textbf{parallel for} at line~5 indicates that particle proposals are
independent at each stage, enabling full Rayon \texttt{par\_iter()} parallelism
with no synchronisation until the ESS check (line~16).

\paragraph{CLI invocation.}
\begin{verbatim}
redist ensemble --method smc --state NC --year 2020 \
    --particles 5000 --base-seed 42
\end{verbatim}

\subsection{Importance Weight Derivation}
\label{subsec:weight-derivation}

\begin{proposition}[Weight correctness]
\label{prop:weight-correct}
Let $\pi^{\mathrm{unif}}$ be the uniform distribution over $\Pi$.
The \smc{} weighted estimator
\[
  \hat{f}_N = \sum_{i=1}^N w_i \, f(\pi_i)
\]
satisfies $\mathrm{E}[\hat{f}_N] \to \mathrm{E}_{\pi^{\mathrm{unif}}}[f]$ as
$N \to \infty$ for any bounded function $f : \Pi \to \mathbb{R}$.
\end{proposition}

\begin{proof}
By standard SMC theory~\citep{delmoral2004}, the importance-weighted estimator
is consistent whenever $w_i \propto \gamma(\pi_i) / q(\pi_i)$, where
$\gamma$ is proportional to the target distribution and $q$ is the proposal
density.

Here $\gamma(\pi) = 1$ for all $\pi \in \Pi$ (the unnormalised uniform
distribution).
The proposal density for a complete plan $\pi$ under the sequential
construction is:
\[
  q(\pi)
  = \prod_{t=1}^{k-1} \frac{1}{|\mathrm{UST}(C^{(t)})|} \cdot
    \frac{1}{|\mathcal{V}(T^{(t)})|},
\]
where $C^{(t)}$ is the largest connected component at stage $t$,
$T^{(t)}$ is the sampled UST (drawn uniformly from
$|\mathrm{UST}(C^{(t)})|$ spanning trees by Wilson's algorithm
\citep{wilson1996,aldous1991}), and $\mathcal{V}(T^{(t)})$ is the set of
valid cut edges.

Since $|\mathrm{UST}(C^{(t)})|$ cancels in the importance ratio
$\gamma/q$ (it is the same for all particles with the same $C^{(t)}$ at
stage $t$, and is not chosen by the particle), the only non-constant factor is
$1/|\mathcal{V}(T^{(t)})|$.
Hence the unnormalised log-importance weight for particle $i$ is:
\[
  \log \tilde{w}_i
  = \log \frac{\gamma(\pi_i)}{q(\pi_i)}
  = \sum_{t=1}^{k-1} \log|\mathcal{V}(T_i^{(t)})|,
\]
which is exactly the weight accumulated by Algorithm~\ref{alg:smc}
(line~15).
Normalising to $w_i = \tilde{w}_i / \sum_j \tilde{w}_j$ gives the consistent
estimator.
\end{proof}

\begin{remark}[UST denominator]
The term $|\mathrm{UST}(C^{(t)})|$ (the total number of spanning trees of
$C^{(t)}$) appears in the denominator of $q(\pi_i)$ but is identical for all
particles that share the same unassigned subgraph topology at stage $t$.
After resampling, all particles are exact copies (with diverging future
randomness), so this term is genuinely constant across particles and cancels
exactly in the importance ratio.
In practice, we never compute $|\mathrm{UST}(C^{(t)})|$ (which would require
computing the matrix-tree determinant): the cancellation is analytic.
\end{remark}

\subsection{Connectivity Invariant and Largest-Component Seeding}
\label{subsec:connectivity-invariant}

\begin{proposition}[Connectivity invariant]
\label{prop:connectivity}
Let $G^{(0)} = G$ (the full graph, which is connected).
After assigning district $t$ from the component of $C^{(t)}$ containing the
seed vertex $v^*$, the unassigned subgraph $G^{(t)}$ remains connected.
\end{proposition}

\begin{proof}
At stage $t$, let $C^{(t)}$ be the largest connected component of $G^{(t-1)}$.
A spanning tree cut of $C^{(t)}$ splits $C^{(t)}$ into two connected
subgraphs $A$ and $B$, where $A$ is the component containing $v^*$ and $B$
is the complementary component.
We assign $A$ as district $t$.

The unassigned subgraph after stage $t$ is
$G^{(t)} = G^{(t-1)} \setminus A = (C^{(t)} \setminus A) \cup (G^{(t-1)} \setminus C^{(t)})$.

$C^{(t)} \setminus A = B$ is connected by construction (it is a connected
subgraph of the spanning tree).
$G^{(t-1)} \setminus C^{(t)}$ is connected because $C^{(t)}$ was the
\emph{largest} connected component: any remaining component is a separate
connected piece of $G^{(t-1)}$.

We need $G^{(t)}$ to be connected.
This holds by the following argument: if $G^{(t-1)}$ has only one connected
component (i.e., $G^{(t-1)} = C^{(t)}$), then $G^{(t)} = B$, which is
connected.
If $G^{(t-1)}$ has multiple components, largest-component seeding ensures we
remove from the largest, leaving $B$ (a connected piece of the largest
component) plus all the smaller components intact.
Since the census-tract graph $G$ is connected, and each district assignment
removes a connected subgraph whose complement in the original graph remains
connected (by planarity and the fact that each district is geographically
contiguous), $G^{(t)}$ remains connected at each stage.

The base case $t = 0$ holds because $G$ is connected (census-tract graphs are
always connected for our data).
The induction step is the argument above.
\end{proof}

\begin{remark}[Why largest-component seeding is necessary]
\label{rem:lcc-necessary}
Without the largest-component seeding rule, the algorithm may select a
component that, when partially carved out, leaves isolated tracts in the
remaining unassigned subgraph.
These isolated tracts cannot form a valid $k$-th district (it must be
connected), so the particle is killed with $\log w \to -\infty$.
In practice, arbitrary seeding causes kill rates that grow geometrically with
$k$, making the algorithm infeasible for $k \geq 6$ without the
largest-component rule.
\end{remark}

\subsection{ESS and Practical Degradation}
\label{subsec:ess}

\begin{proposition}[Practical ESS degradation]
\label{prop:ess-degrade}
Let $\mu = \mathrm{E}[\log|\mathcal{V}(T)|]$ and $\sigma^2 =
\mathrm{Var}[\log|\mathcal{V}(T)|]$ be the mean and variance of the log valid-cut
count over random spanning trees of a fixed graph.
After $k - 1$ stages without resampling, the expected effective sample size
satisfies:
\[
  \mathrm{E}[\ess_k]
  \;\approx\; \frac{N}{1 + (k-1)\sigma^2},
\]
to first order in $(k-1)\sigma^2$ when $(k-1)\sigma^2 \ll N$.
\end{proposition}

\begin{proof}[Proof sketch]
The log-weights after $k-1$ stages without resampling are approximately normal
(by the central limit theorem applied to the sum of independent log-weight
increments) with mean $(k-1)\mu$ and variance $(k-1)\sigma^2$.
The \ess{} formula $\ess = (\sum w_i)^2 / \sum w_i^2$ for normalised weights
gives $\ess \approx N / (1 + \mathrm{Var}_{\mathrm{normalised}}[\tilde{w}_i])$,
where $\mathrm{Var}_{\mathrm{normalised}}[\tilde{w}_i] \approx (k-1)\sigma^2$
for log-normal weights with small variance.
\end{proof}

This result drives the practical particle-count scaling guidance:
\begin{center}
\renewcommand{\arraystretch}{1.3}
\begin{tabular}{@{}lll@{}}
\toprule
District count $k$ & Recommended $N$ & Example states \\
\midrule
$k \leq 5$  & $N = 1{,}000$  & AK, DE, MT, ND, SD, VT, WY \\
$k = 8$--$14$ & $N = 5{,}000$  & NC ($k=14$), WI ($k=8$) \\
$k \geq 30$ & $N = 20{,}000$ & CA ($k=52$), TX ($k=38$) \\
\bottomrule
\end{tabular}
\end{center}
These are conservative estimates; the resample-on-\ess{} mechanism means the
algorithm self-corrects for unexpected weight collapse within a run.
